\renewcommand{\IEEEkeywordsname}{Palavras-chave}

\begin{abstract}
Este relatório apresenta a primeira parte do projeto de Ciência de Dados sobre os empréstimos da LendingClub (2007--2018): formulação do problema, análise e limpeza dos dados e seleção de modelos. As principais tarefas foram prever a taxa de juro antes de esta ser fixada (regressão), prever na concessão se o empréstimo vai entrar em \textit{default} (classificação) e agrupar mutuários e estados (clustering) e, para cada uma, definir o momento de previsão, as variáveis elegíveis, a coorte e uma divisão temporal com teste em 2015. A análise exploratória, feita apenas com dados de desenvolvimento, mostra que a taxa depende sobretudo da qualidade do crédito e do prazo, que o grau atribuído pela plataforma separa bem o risco (5,9\% de \textit{default} no grau A contra 35,9\% no grau G) e que as relações das restantes variáveis com o \textit{default} são fracas e, muitas vezes, não lineares. Com base nestes resultados, criam-se dez \textit{features} e define-se um plano de modelação com kNN, modelos lineares regularizados, árvores de decisão, \textit{Random Forest}, LightGBM e redes neuronais, validados com uma janela temporal expansível e comparados com \textit{baselines} simples.
\end{abstract}